\chapter{绪论}

\section{研究背景与意义}

在量化投资研究中，股票收益并不是由单一指标决定的，而是由估值、成长、质量、动量、波动率、流动性以及宏观环境等多种因素共同驱动。随着数据采集技术和量化平台的发展，研究者能够同时获得财务报表数据、日频交易数据、技术指标、行业分类、公告事件以及指数成分变动等多源信息，传统的少量核心因子分析范式逐渐演化为高维多因子分析范式。对于单只股票而言，其在不同交易时点上往往对应几十甚至上百个候选因子；对于整个股票池而言，数据天然包含股票截面、因子集合和时间演化三个维度。因此，股票因子研究的基础对象已经不再是简单的二维矩阵，而是具有显著多维耦合特征的高阶数据结构。

过去较为常见的做法是先将高维因子数据展平为“样本$\times$特征”的矩阵，再使用主成分分析、逐步回归、相关系数筛选或聚类方法进行降维。这类方法的优点是计算简单、工程成熟，但其隐含假设是样本之间相互独立、特征之间关系稳定，且时间维度只作为额外标签存在。对于股票因子问题而言，上述假设往往并不充分成立：同一类风格因子在不同市场阶段的表现可能显著变化，不同行业股票对同一因子的敏感度也并不一致，因而仅用二维矩阵难以完整表达“股票—因子—时间”的耦合结构。若在建模初期就将高维结构展平，往往会损失跨模态关联信息，最终使降维结果在解释性与稳定性上受到限制。

张量分解方法为这一问题提供了更自然的建模路径。张量是矩阵的高阶推广，能够在保持原始多维结构的前提下进行低维表示学习。其中，CP 分解通过若干秩一张量之和表示原始张量，适合提取共享的潜在成分；Tucker 分解通过核张量与多个因子矩阵的乘积刻画不同模态间更复杂的交互结构，具有较强的灵活性。对于股票因子研究来说，张量分解不仅可以完成高维因子压缩，还能够同时输出股票模态、因子模态和时间模态的潜在表示，从而为股票联动模式识别、因子共振关系发现以及时间状态切换分析提供统一基础。

从理论意义上看，本文将张量分解应用于股票因子降维与模式发现，能够在金融数据分析语境下验证高阶张量建模的适用性，并进一步比较 CP、Tucker 与 PCA 三种路径在结构保持、信息提取和结果解释上的差异边界。一方面，这有助于将张量分解的理论框架从信号处理、计算机视觉等领域延伸到量化金融场景；另一方面，也能够为“多因子体系如何在不牺牲结构信息的前提下实现有效压缩”这一问题提供更系统的方法论支持。

从实践意义上看，若能在高维股票因子集合中识别出少量信息密度更高、结构解释更清晰的低维核心表示，就可以降低选股模型的复杂度，减轻多重共线性问题，提升后续风格解释和风险监控的稳定性。此外，基于张量分解输出的潜在时间表示，还可以用于观察市场状态切换、风格轮动和异常时期的结构突变，这对构建更稳健的量化研究流程和实证分析框架具有明显现实价值。

\section{国内外研究现状}

张量分解的数学思想可以追溯到 Hitchcock 对多元数组表示形式的研究。Hitchcock 在 1927 年给出了将高阶张量表示为多个乘积项求和的基本形式，为后续规范分解奠定了数学起点\textsuperscript{[1]}。随后，Tucker 在三模因子分析中提出核张量结构，使高阶数据可以通过不同模态的潜在因子矩阵进行表达\textsuperscript{[2]}；Harshman 又从多模态因子分析角度发展了 PARAFAC 方法，使 CP 分解逐步成为高阶结构分析中的经典模型\textsuperscript{[3]}。Kolda 与 Bader 的综述系统总结了张量分解的数学定义、算法框架和应用场景，被广泛视为该领域的重要基础文献\textsuperscript{[4]}。胡胜龙关于张量分解唯一性的研究则进一步梳理了 CP 分解的可辨识条件和理论边界，为高维张量模型的解释性提供了更扎实的理论支撑\textsuperscript{[5]}。

从方法综述角度看，夏虹对张量方法及其应用的综述说明，张量建模的最大优势在于能够直接处理多维结构数据，而不需要在建模初期将多模态信息粗糙展平\textsuperscript{[6]}。Rabanser、Shchur 与 Günnemann 的综述则从机器学习视角系统介绍了 CP、Tucker 及其在主题模型、关系数据和时序问题中的应用，指出张量分解已从纯数学工具演化为高维数据分析中的通用方法族\textsuperscript{[7]}。这些研究共同表明，张量方法的核心价值并不仅仅在于“高维”，更在于“结构保持”。

在金融数据分析方向，已有研究逐步将张量方法用于股票波动、金融时序预测和高维参数压缩。岳欣将股票多指标时间序列表示为三阶张量，并利用张量补全方法处理停牌和缺失问题，说明张量方法可以更自然地适应金融数据中的结构性缺失\textsuperscript{[8]}。马少沛等从统计学习角度提出张量充分降维框架，指出若响应变量只依赖张量在低维子空间上的投影，则张量表示有助于兼顾降维效率与信息保留\textsuperscript{[9]}。曾亚丽围绕金融时间序列中的多维耦合结构展开研究，强调张量模型能够把资产、指标与时间的联合变化纳入统一建模框架\textsuperscript{[10]}。Brandi 等关于股票相关结构的研究则指出，仅依赖二维收益相关系数不足以解释股票之间复杂关系，张量化分解能够从更高维角度揭示股票间的共振与联动\textsuperscript{[11]}。关于市场可预测性和时序系统的研究也表明，张量分解在金融市场结构提取、时序参数压缩和潜在因子预测上具有明显潜力\textsuperscript{[12][17]}。

在多因子选股领域，传统路径仍以回归分析和统计筛选为主。黄宏运等基于多元回归构建多因子选股模型，从基本面和技术面因子中进行组合筛选，体现了传统统计方法在多因子建模中的可操作性\textsuperscript{[13]}。刘宇轩则在多因子模型中引入金融周期指标，对选股模型进行扩展和修正，说明因子体系会随着市场状态和宏观周期变化而发生结构调整\textsuperscript{[14]}。与此同时，高维因子模型研究与张量因子秩选择研究指出，如何在保留有效信息的同时控制潜在维度，是高维金融建模能否成功的关键问题\textsuperscript{[15][16]}。这些成果共同说明，多因子研究已经从“少量因子筛选”转向“高维结构压缩与解释并重”的阶段，而这正是本文希望回应的问题。

综合来看，现有研究已经证明了张量分解在高维结构数据上的理论优势，也展示了其在金融时序、股票波动和相关结构建模中的应用潜力，但仍存在三个值得继续深入的问题。第一，专门针对“股票因子—多维结构保持降维—模式发现”的一体化研究仍不充分，尤其是针对中国 A 股正式样本的系统实验较少。第二，许多研究偏重于预测精度，较少同时讨论因子解释、股票联动和时间状态切换三类结果。第三，工程实现与学术叙事常常彼此脱节，实验产物、数据链路和最终论文结论之间缺乏统一证据来源。本文的研究正是围绕这些不足展开，希望将张量分解方法、正式数据链与结构化实验输出统一到同一研究框架中。

\section{研究内容与技术路线}

本文围绕“股票—因子—时间”三维张量构建正式研究对象，将研究内容划分为四个层面。第一，在数据层，构建可复用的 A 股 formal 数据主链，完成 K 线、财务、报告和指数成分历史等基础数据的整理，并在此基础上形成因子面板。第二，在方法层，比较 CP 分解、Tucker 分解与 PCA 三种降维路径，重点分析张量方法在结构保持和潜在模式提取上的优势。第三，在结果层，围绕分解质量、排序有效性、时间状态变化和因子重要性展开实证分析。第四，在工程层，将实验输出、论文表格、图形可视化和交付材料统一到可追溯的同一仓库中，保证研究证据的一致性。

对应的技术路线可以概括为“原始数据获取与整理—因子面板构建—三阶张量表示—张量分解与秩选择—潜在表示分析—实验评估与可视化输出”。其中，原始数据整理并不是附属步骤，而是后续张量建模能否成立的前提。由于股票市场数据天然存在停牌、成分股变化、财务披露时滞和不同数据源更新口径不一致等问题，本文在正式实验之前，先对 K 线前复权口径、财务/报告覆盖、指数成分历史和因子面板时间窗口进行统一，确保模型输入具有一致的时间跨度和样本边界。

在算法路线方面，本文将先构造原始张量，再执行协议化预处理和训练/验证/测试切分，并分别在 CP、Tucker 和 PCA 三个模型上进行 rank 选择和 held-out 评估。与仅比较重构误差的实验不同，本文希望同时回答三个问题：其一，哪种方法在重构质量上更优；其二，哪种方法在排序指标和滚动稳定性上表现更平衡；其三，哪种方法更有利于形成可解释的因子和时间模式。围绕这三个问题，最终正文将给出表格、图形和定性讨论。

\section{论文结构安排}

除摘要、参考文献、附录等部分外，全文共分为六章。第一章为绪论，介绍研究背景、国内外研究现状、研究内容与技术路线。第二章阐述张量代数、张量分解与金融因子数据表示的理论基础。第三章给出数据来源、样本边界、预处理流程和张量构造方式。第四章详细说明基于张量分解的股票因子降维与模式发现方法。第五章基于正式实验结果，围绕分解质量、模式识别与排序有效性展开分析。第六章总结全文研究结论、创新点、不足与后续展望。

\chapter{理论基础与关键技术}

\section{张量代数基础}

在本文中，股票因子数据被表示为三阶张量$\mathcal{X}\in\mathbb{R}^{N\times F\times T}$，其中$N$表示股票数量，$F$表示因子数量，$T$表示观测时点数。与向量和矩阵相比，张量最大的特点在于能够同时保留多个模态的信息，因此更适合描述高维结构数据。若将张量简单展平成矩阵，则本属于不同模态的结构特征会被强行混合，造成解释上的模糊。基于此，本文直接在张量形式下进行建模。

对三阶张量而言，最常用的基础运算包括模态展开、外积、Kronecker 积、Khatri--Rao 积和 Hadamard 积。设$\mathcal{X}$在第一模上的展开矩阵为$X_{(1)}$，则其维度为$N\times(FT)$；类似地，第二模和第三模展开分别记为$X_{(2)}$和$X_{(3)}$。模态展开的意义在于，它能够将高阶张量上的某一步优化转化为矩阵上的最小二乘问题，从而为交替最小二乘等算法提供计算基础。若记向量$a,b,c$的外积为$a\circ b\circ c$，则其结果是一个秩一三阶张量，任意元素满足
\begin{equation}
\left(a\circ b\circ c\right)_{ijt}=a_i b_j c_t.
\end{equation}
高阶张量的许多分解实际上都可以理解为若干秩一张量的叠加或由低维因子矩阵在不同模态上作用得到的结构重组。

在股票因子应用中，张量代数的价值不仅体现在数学表示上，也体现在金融解释上。股票模态对应横截面结构，因子模态对应信息维度，时间模态对应市场演化。若某一潜在成分在股票模态上的载荷主要集中于银行、券商和保险股，而在时间模态上又集中出现在政策宽松阶段，则该潜在成分就可能对应“金融风格因子在特定市场环境下的整体共振”。因此，张量的每个模态都具有明确的经济含义，而不是单纯的索引维度。

\section{张量分解模型}

本文重点使用 CP 分解与 Tucker 分解两类张量方法。CP 分解的核心思想是将原始张量表示为若干秩一张量之和，即
\begin{equation}
\mathcal{X}\approx \sum_{r=1}^{R}\lambda_r a_r\circ b_r\circ c_r.
\end{equation}
其中，$R$是分解秩，$\lambda_r$为第$r$个成分的权重，$A=[a_1,\ldots,a_R]$、$B=[b_1,\ldots,b_R]$、$C=[c_1,\ldots,c_R]$分别对应股票、因子和时间三个模态的潜在因子矩阵。对股票因子问题而言，矩阵$B$最直接反映了原始因子对潜在成分的贡献强度，矩阵$A$可用于分析股票在低维潜在空间中的相似性，矩阵$C$则可用于观察市场状态的时间变化。

CP 分解的优点在于结构简洁、参数量相对可控，且当唯一性条件满足时，分解结果具有较强可解释性\textsuperscript{[4][5]}。然而，CP 分解也面临秩选择困难和局部最优敏感的问题。尤其在金融数据场景下，若不同模态之间的交互关系较为复杂，单纯依赖秩一张量叠加可能不足以准确刻画全部结构，这就需要引入更灵活的 Tucker 分解。

Tucker 分解将三阶张量表示为核张量和多个因子矩阵的多模态乘积：
\begin{equation}
\mathcal{X}\approx \mathcal{G}\times_1 A\times_2 B\times_3 C.
\end{equation}
其中，$\mathcal{G}$是核张量，用于刻画不同潜在成分之间的高阶交互关系；$A$、$B$、$C$分别是股票、因子和时间模态上的因子矩阵。与 CP 分解相比，Tucker 分解允许不同模态设置不同秩，例如在股票模态保留三个潜在维度、因子模态保留两个潜在维度、时间模态保留两个潜在维度，这使得模型更适合描述非对称的复杂结构。

在算法求解上，本文主要依赖交替最小二乘和基于高阶奇异值分解的初始化策略。以 CP 分解为例，当固定$B$和$C$后，对$A$的更新可以写为
\begin{equation}
A\leftarrow X_{(1)}(C\odot B)\left[(C^{\top}C)\ast(B^{\top}B)\right]^{-1},
\end{equation}
其中$\odot$表示 Khatri--Rao 积，$\ast$表示 Hadamard 积。对于 Tucker 分解，则可以先采用 HOSVD 或截断奇异值分解获得初始因子矩阵，再通过交替迭代优化核张量和各模态因子矩阵。总体而言，CP 分解对应结构更为简洁的低秩张量表达，而 Tucker 分解则能提供更灵活的多线性秩约束，在重构能力与结构表达之间形成更平衡的结果。

\section{金融数据的张量表示方法}

在股票因子问题中，最自然的张量表示方式是以股票为第一模、因子为第二模、时间为第三模，将每个元素$x_{ijt}$定义为第$i$只股票在第$t$个交易日上的第$j$个因子取值，即
\begin{equation}
\mathcal{X}=(x_{ijt})\in\mathbb{R}^{N\times F\times T}.
\end{equation}
这一表示方式具有三方面优势。第一，它保留了股票横截面之间的相对关系，可直接观察行业、风格和规模层面的结构接近性。第二，它保留了因子之间的共振关系，适合识别冗余因子和具有替代关系的因子簇。第三，它保留了时间方向上的连续性和阶段性，能够描述因子结构在牛熊转换、流动性冲击和政策扰动中的演化。

若进一步考虑正式实验中的样本池边界，则同样可以在统一结构下表示 HS300、SZ50 与 ZZ500 三个样本池对应的因子张量，只需在样本选择阶段通过指数成分历史对股票集合进行限定。由于本文目标并非构建一个全市场实时交易系统，而是完成结构保持降维与模式发现，因此更关注样本口径一致、时间窗口完整和特征构造稳定，而不是在不同样本池之间混合建模。

\section{模型评价指标}

为了评价降维效果，本文使用重构误差、解释方差、排序相关、信息比率和滚动稳定性等指标。首先，重构误差衡量分解后张量与原始张量之间的偏离程度，可写为
\begin{equation}
\mathrm{MSE}=\frac{1}{NFT}\|\mathcal{X}-\hat{\mathcal{X}}\|_F^2.
\end{equation}
与之对应，解释方差反映模型对原始张量波动信息的保留程度，数值越大通常表示降维后保留的有效结构越多。

其次，考虑到本文不仅关注重构能力，也关注低维表示的排序价值，因此还使用 Rank IC 和 IR 指标来观察降维后的潜在得分与下一期收益之间的相关关系。其中，Rank IC 侧重于秩相关，适合金融场景下强调排序而非绝对值预测的任务。若某模型在解释方差上较高但 Rank IC 很低，则意味着它更擅长数值重构而不一定更适合选股排序。

最后，本文引入滚动稳定性指标衡量模型在不同时间窗口下输出潜在结构的一致性。对因子降维而言，若模型在不同窗口中频繁产生完全不同的潜在结构，则即使单次拟合表现良好，也很难形成可复用的研究结论。因此，重构误差、解释方差、Rank IC 与滚动稳定性共同构成本文的评价体系，其中前两者侧重数值拟合，后两者侧重金融解释与研究稳健性。

\chapter{数据预处理与张量表示}

\section{数据来源与样本说明}

本文使用的数据来源于项目内部构建的 formal 数据主链。该链路以 Baostock 为主要原始数据接口，围绕日频 K 线、财务数据、业绩报告、指数成分历史、分红复权信息和衍生因子面板形成统一的数据组织方式。正式样本池分别为 HS300、SZ50 和 ZZ500。需要说明的是，配置文件请求的样本窗口为 2015 年 1 月 1 日至 2026 年 4 月 1 日，但当前根仓库中已提交的正式实验产物实际覆盖窗口为 2026 年 3 月 2 日至 2026 年 3 月 30 日，这一点在运行 manifest 中可以直接核对。之所以仍以这三个样本池为主，是因为它们具有较清晰的风格边界、较好的流动性基础和较稳定的成分定义，能够更清楚地观察因子结构与样本池风格之间的对应关系。

从数据覆盖角度看，当前正式工作链路已经完成前复权 K 线、财务表、报告表和指数成分历史的整理，但这些原始层数据与当前根仓库中真正提交的实验产物之间仍然存在窗口差异。也就是说，数据底座具备继续向长期窗口扩展的基础，而当前已提交的 formal 因子面板与正式运行结果仍主要体现为 2026 年 3 月的短窗口实验版本。这也是本文在局限性分析中专门强调“正式窗口口径仍需进一步统一”的原因。

需要强调的是，“数据已经抓取完成”与“数据已经安全进入训练输入”并不是同一概念。当前正式训练张量仍以主线因子面板为核心，但已经在 extended 配置中接入了第一版市场级代理变量、财务 PIT 指标以及业绩快报/业绩预告事件特征；更完整的宏观变量和更广泛的事件字典则尚未系统并入统一训练口径。为避免在训练阶段引入披露滞后和未来信息泄露，本文另外形成了《训练输入扩展与 PIT 安全说明》，将扩展特征划分为市场代理、财务点时特征和事件特征三类，并分别给出特征字典建立、披露时点映射、可用交易日判断和分层扩展实验顺序。该说明既是当前训练输入边界的补充解释，也是后续扩展实验的执行依据。

在样本组织上，本文不直接把所有股票合并到同一张量中统一训练，而是分别对 HS300、SZ50 和 ZZ500 构建因子张量。这种做法的动机在于，不同指数样本池本身具有明显不同的规模特征、行业构成和风格偏向。若直接混合，潜在表示可能首先学习到“样本池差异”，而非股票因子内部的结构关系。分样本池建模则更有利于在同质性更高的截面内观察张量分解的表现。

\section{股票池与时间窗口设定}

指数样本池的边界并非静态固定，因此本文采用“指数成分历史”而非单一时点快照定义股票池。对于任一交易日，只将当日属于目标指数成分集合的股票纳入样本。这样做有两个好处：第一，避免了使用当前成分股去回溯历史样本而产生的未来信息泄露；第二，能够更真实地反映指数样本池在长期窗口中的构成变化。

从研究设计角度看，更理想的正式窗口应覆盖 2015 年 1 月至 2026 年 4 月，以便同时观察多轮市场阶段变化；但就当前根仓库已提交的正式实验结果而言，实际落盘窗口仍是 2026 年 3 月。因此，本文后文中涉及的实证结论应理解为“当前提交版实验产物对应窗口”下的结果，而不是已经完全落到仓库中的长期全窗口结论。

\section{数据预处理}

金融因子数据在进入张量建模之前，必须经过严格预处理。首先是缺失值处理。股票因子缺失往往来自停牌、上市时间不足、财务数据尚未披露或数据接口限制等原因。若直接删除含缺失值样本，会在长期窗口下造成大量信息浪费，并使不同样本池之间出现不可控的不平衡。因此，本文在项目实现中采用“训练集拟合预处理状态，再统一作用于验证集与测试集”的协议化流程，并结合前向填充、后向填充及行业中位数填充减少缺失影响。

其次是异常值处理。A 股市场中部分因子，尤其是估值类和收益率类因子，容易受到极端行情、个股停复牌或低基数效应影响。若不处理异常值，张量分解得到的潜在表示会过度受极端点主导，降低整体稳定性。为此，本文采用截尾或中位数绝对偏差等方式对极端值进行收缩，再执行标准化。标准化的目的并不是简单地把所有因子拉到同一数值区间，而是为了使不同因子在分解时具备可比较性，避免某一类量纲较大的因子在 Frobenius 范数意义下占据过高权重。

再次是时间与样本切分协议。本文没有采用“整段样本一次性拟合再解释”的简单流程，而是显式划分训练、验证和测试分区。预处理状态只在训练集上拟合，验证集用于选择 rank，测试集只用于最终 held-out 评估。这样的流程使论文中的实验结果不至于因为预处理泄露而过度乐观，也更贴近正式机器学习评估规范。

最后是辅助数据与主线输入的边界控制。虽然当前 formal 数据链已经具备财务、报告、复权和更多市场级数据资产，本文在 baseline 训练输入中仍以已经稳定进入主线的因子面板为核心；而在 extended 配置中，仅把已经形成 PIT 合同和可交易日规则的市场代理变量、财务 PIT 与事件特征纳入训练，其余仍处于补充合同阶段的数据继续保留在查询和解释层。这样做的目的，是保证实验口径明确、建模边界清晰，避免因补充数据口径尚未完全稳定而影响主体结论。

\section{张量表示}

在完成预处理之后，本文将每个样本池的数据表示为三阶张量$\mathcal{X}\in\mathbb{R}^{N\times F\times T}$。其中，$N$是该样本池在给定交易日有效成分股集合下的股票数量，$F$是因子面板中的因子数，$T$是时间维度上的交易日数。由于指数成分历史会导致股票集合在不同时间上变化，工程实现中采用了统一编码与缺失补位的方式，以保证张量在模态维度上保持固定形状。

这种张量表示与传统矩阵表示的最大区别在于，时间维度不再只是“样本编号的一部分”，而是与股票和因子维度共同组成潜在结构的核心模态。由此得到的模型结果不仅包含重构误差等数值信息，也会输出三个模态上的潜在表示矩阵，从而为后续的股票相似性、因子重要性和时间状态变化分析提供直接依据。

\chapter{基于张量分解的股票因子降维与模式发现方法}

\section{问题定义}

本文研究的问题可以形式化为：给定股票—因子—时间三阶张量$\mathcal{X}$，如何在保持其多维结构信息的前提下，构造一个低维表示$\hat{\mathcal{X}}$，使得其既能较好重构原始张量，又能输出对股票、因子与时间具有解释意义的潜在表示。相比一般的“仅降维”问题，本文更强调“降维+模式发现”的双重目标。

具体而言，本文希望完成以下三个子任务。第一，在因子模态上实现有效压缩，从而降低高维因子集合中的冗余程度。第二，在股票模态上识别潜在的相似结构，用于解释哪些股票在潜在风格暴露上更接近。第三，在时间模态上识别市场阶段变化，观察潜在结构在长期样本中的迁移与突变。围绕这三个目标，单纯使用二维 PCA 虽然可以提供重构基线，但不具备自然的多模态解释能力，因此需要引入张量分解方法。

\section{张量分解模型构建}

本文同时构建 CP、Tucker 与 PCA 三类模型。CP 模型作为最简洁的张量分解基线，用于观察在低秩约束较强时，三阶张量能否被少量公共成分有效表示。Tucker 模型作为论文主方法，用于在不同模态设置不同秩，捕捉更复杂的模态间交互结构。PCA 模型虽然不在张量形式下建模，但由于其在传统降维问题中具有较强数值表现，因此作为二维压缩基线保留。

对 CP 分解而言，模型秩$R$越大，理论上可表示的结构越丰富，但解释性和压缩率会下降；秩过小则可能无法表示足够的结构信息。对 Tucker 分解而言，需要同时设置股票模态秩、因子模态秩和时间模态秩。本文在实验中并不依赖人工经验一次性固定秩，而是先给出候选 rank 集合，再在验证集上比较不同 rank 对应的 held-out MSE，选出表现最优的参数组合。这种做法能在一定程度上避免“为了论文结论而人为挑选 rank”的主观性问题。

在因子筛选与模式解释上，本文主要依赖因子模态矩阵$B$及其载荷大小。若某个潜在成分在$B$中对质量因子、估值因子和动量因子分别给出不同权重，则说明不同金融含义的因子被组合进不同潜在方向中。进一步地，可定义因子$j$的综合重要性为
\begin{equation}
w_j=\frac{\|b_{j:}\|_2}{\sum_{k=1}^{F}\|b_{k:}\|_2},
\end{equation}
并按$w_j$排序识别核心因子。这一做法虽然不直接等价于经济学中的因子定价显著性检验，但对于“哪些因子在低维潜在结构中承担了主要信息功能”这一问题具有直观解释力。

\section{模型求解与实现}

工程实现上，本文的实验流程可以概括为五步。第一，读取 formal 样本池的因子面板和指数历史，构造原始三阶张量。第二，按照显式训练/验证/测试协议执行预处理和切分。第三，对每一类模型分别在候选 rank 集合上进行训练和验证，选择最优参数。第四，在训练集与验证集联合数据上重拟合最终模型，并只在测试集上进行 held-out 评估。第五，将所有结果写出为结构化产物，包括 metrics、factor summary、time regimes、selection candidates 和图形可视化。

在算法层面，CP 和 Tucker 模型均依赖迭代更新。为了减轻局部最优对结果的影响，实验实现中使用固定随机种子和稳定的预处理协议，并在训练日志中记录 rank 选择过程和最终输出路径。PCA 则作为二维基线，以相同测试集口径进行评估，确保三类方法比较公平。需要指出的是，本文关注的是研究型实证框架，而不是生产级高频交易系统，因此算法实现更强调可复现性、可追溯性和结构化结果输出，而非极限运行速度。

\chapter{实验结果与分析}

\section{实验设置}

正式实验样本池包括 HS300、SZ50 与 ZZ500。当前根仓库中可直接复现的 formal profile、factor panel 快照与输出目录已统一到 2026 年 3 月 2 日至 2026 年 3 月 30 日窗口。每个样本池均基于正式因子面板构建张量，并在统一的预处理、切分和评价协议下比较 CP、Tucker 与 PCA 三种方法。其中，CP 和 Tucker 负责提供张量结构建模结果，PCA 作为二维压缩基线用于比较其数值拟合能力。

从输出结果看，实验不仅生成数值指标，还生成了模型解释方差对比图、Rank IC 对比图、模型指标总览图、时间状态变化图、因子重要性热力图，以及组合累计净值对比图和回撤曲线对比图。这些图现已作为结构化 SVG 产物写入 formal 输出目录，可与表格和 JSON 结果一一对应。这样的图文并行输出方式有助于把论文从“结果表格堆砌”转向“结构模式分析”，也是本文更强调的一点。

\section{分解质量分析}

首先从数值重构能力出发，三组样本池上的结果具有较清晰的一致性。HS300 样本中，CP、Tucker、PCA 的 MSE 分别为 0.9140、0.3406 和接近 0，对应解释方差分别为 0.0017、0.6280 和 1.0000；SZ50 样本中，三者的 MSE 分别为 0.9898、0.3160 和接近 0，对应解释方差分别为 $-0.0527$、0.6639 和 1.0000；ZZ500 样本中，三者的 MSE 分别为 1.0091、0.4832 和接近 0，对应解释方差分别为 $-0.0541$、0.4953 和 1.0000。可以看到，在当前短窗口 formal 输出下，PCA 在三组样本池上都表现出最强的重构能力，而 Tucker 在两个张量方法中稳定优于 CP。

这一结果说明，若只以数值重构为目标，二维压缩方法仍具有明显优势。原因在于，PCA 在矩阵展开后直接针对方差最大方向建模，天然更偏向“数值逼近最优”；而张量方法在保持多模态结构约束的同时，必须在不同模态间平衡表示能力，因此纯重构性能未必领先。然而，这并不意味着张量方法失去研究价值。恰恰相反，若研究目标是结构解释与模式发现，那么在重构精度略有让步的情况下换取更清晰的多模态潜在表示，是完全合理的策略。

\begin{figure}[!t]
\centering
\begin{tikzpicture}
\begin{groupplot}[
  group style={group size=3 by 1, horizontal sep=1.3cm},
  width=0.28\textwidth,
  height=6.0cm,
  ybar,
  ymin=-0.1,
  ymax=1.0,
  symbolic x coords={CP,Tucker,PCA},
  xtick=data,
  nodes near coords,
  nodes near coords align={vertical},
  every node near coord/.append style={font=\tiny},
  tick label style={font=\small},
  ylabel style={font=\small},
  title style={font=\small},
  legend style={font=\small, draw=none}
]
\nextgroupplot[title={HS300},ylabel={解释方差}]
\addplot[fill=blue!55] coordinates {(CP,0.0017) (Tucker,0.6280) (PCA,1.0000)};
\nextgroupplot[title={SZ50}]
\addplot[fill=orange!70] coordinates {(CP,-0.0527) (Tucker,0.6639) (PCA,1.0000)};
\nextgroupplot[title={ZZ500}]
\addplot[fill=teal!65] coordinates {(CP,-0.0541) (Tucker,0.4953) (PCA,1.0000)};
\end{groupplot}
\end{tikzpicture}
\caption{三个正式样本池上的模型解释方差对比图}
\label{fig:variance}
\end{figure}

图~\ref{fig:variance} 更直观地展示了解释方差差异。PCA 在三个样本池上均显著领先；Tucker 稳定保持较高的正解释方差，而 CP 只有在 HS300 上勉强保持接近于零的解释方差，在 SZ50 和 ZZ500 上仍为负值。这表明，在当前根仓库提交的 formal 窗口内，低秩 CP 仍不足以稳定表示股票因子结构，而 Tucker 在更灵活的多线性秩约束下能够更稳定地提取结构信息。

\section{模式发现结果分析}

除了数值重构能力外，本文更关注潜在表示所揭示的模式结构。根据因子重要性输出，价值因子、质量因子和波动率因子在多个样本池和多个模型下都反复出现于前列，说明它们构成了当前正式因子面板中的核心信息轴。以当前 HS300 的 Tucker 结果为例，波动率因子、价值因子和质量因子的平均重要性依次居前，而动量因子在当前短窗口内的重要性接近于零，说明在这一窗口下市场结构更偏向估值、质量与波动信息，而非长周期趋势信息。

在股票模态层面，潜在表示可用于构造股票相似关系。虽然本文不把股票相似矩阵单独作为最终评价指标，但从实验日志和相似对输出可以观察到，同一行业或相近风格股票更容易在潜在空间中表现出接近的载荷结构。这与量化投资中“因子暴露决定风格分组”的直觉是一致的。相较于简单相关系数，张量潜在表示构造的相似关系同时考虑了时间和因子两个方向，因此更容易反映结构性接近，而非仅仅是短期价格同步。

本轮进一步基于全 A 样本的 Tucker 选择结果生成了股票潜在结构图与 cluster-vs-industry 交叉图。对应产物位于 \texttt{code/data/formal/reports/pattern\_discovery/}，其中 \texttt{stock\_latent\_structure\_formal\_all\_a\_run\_tucker.svg} 将股票按潜在聚类标签展开，并用行业颜色展示聚类内部的行业构成；\texttt{cluster\_vs\_industry\_formal\_all\_a\_run\_tucker.svg} 则把聚类标签与行业分布交叉统计。该图组说明，股票潜在空间并不是简单复制行业分类，而是在行业信息之外保留了由因子暴露、时间状态和排序得分共同形成的结构相似性。

时间模态则提供了对市场状态切换的观察窗口。正式实验输出的 \texttt{time\_regimes\_tucker.json} 记录了若干显著的相邻日期结构跳变。在当前已提交的 HS300 正式结果中，最明显的跳变集中在 2026 年 3 月 25 日至 2026 年 3 月 30 日这一短区间，说明即使在当前较短的正式输出窗口内，潜在结构也并非静态稳定，而是在局部阶段出现明显重排。换言之，张量分解不仅用于“压缩因子”，还能够识别“市场何时发生了潜在结构变化”。

\begin{figure}[!t]
\centering
\begin{tikzpicture}
\begin{groupplot}[
  group style={group size=3 by 1, horizontal sep=1.3cm},
  width=0.28\textwidth,
  height=6.0cm,
  ybar,
  ymin=-0.03,
  ymax=0.05,
  symbolic x coords={CP,Tucker,PCA},
  xtick=data,
  nodes near coords,
  nodes near coords align={vertical},
  every node near coord/.append style={font=\tiny},
  tick label style={font=\small},
  ylabel style={font=\small},
  title style={font=\small}
]
\nextgroupplot[title={HS300},ylabel={Rank IC}]
\addplot[fill=blue!55] coordinates {(CP,0.0833) (Tucker,0.0731) (PCA,0.0174)};
\nextgroupplot[title={SZ50}]
\addplot[fill=orange!70] coordinates {(CP,-0.0167) (Tucker,0.1523) (PCA,0.0129)};
\nextgroupplot[title={ZZ500}]
\addplot[fill=teal!65] coordinates {(CP,-0.0779) (Tucker,0.0573) (PCA,-0.0178)};
\end{groupplot}
\end{tikzpicture}
\caption{三个正式样本池上的模型 Rank IC 对比图}
\label{fig:rankic}
\end{figure}

图~\ref{fig:rankic} 展示了三个正式样本池在 Rank IC 指标上的差异结构。可以看到，Tucker 在三个样本池上都保持了正向表现，而 CP 只在 HS300 上取得了 0.0833 的正向 Rank IC；PCA 则在 HS300 与 SZ50 上保留了较弱的正向排序相关。这意味着仅有较强重构能力并不足以保证排序有效性，保持三维结构的潜在表示在排序任务中仍然具有独立价值。

\section{选股有效性分析}

为了从金融意义上评价降维结果，本文进一步考察 Rank IC 指标。图~\ref{fig:rankic} 显示，在 HS300 样本中，CP 与 Tucker 的 Rank IC 均值分别为 0.0833 和 0.0731，均高于 PCA 的 0.0174；在 SZ50 样本中，Tucker 为 0.1523，PCA 为 0.0129，而 CP 为 $-0.0167$；在 ZZ500 样本中，只有 Tucker 维持了 0.0573 的正向排序相关，而 CP 与 PCA 均为负值。可见，若从“跨样本池保持正向排序相关”这一视角出发，Tucker 在当前正式窗口内仍然表现最为均衡。

这一结果说明，数值拟合最强的方法不一定也是排序最有效的方法。PCA 在解释方差上显著领先，但在 ZZ500 上的 Rank IC 仍为负值，意味着其压缩方向更偏向静态方差解释，而不一定适合直接服务于股票横截面排序。相反，Tucker 通过保持三维结构，在因子压缩的同时保留了更多与时间演化和股票截面相关的模式，因此在三个样本池上都没有出现 Rank IC 由正转负的情况。

从选股研究角度看，这一结论具有现实含义。实际量化研究并不一定要求低维表示精确重构每一个原始因子值，而是更关注压缩后的潜在表示能否支持稳健的股票排序与风格解释。若一个模型在重构误差上最优、但排序效果不稳定，则它更接近数值压缩工具；若一个模型在排序指标和滚动稳定性上都更均衡，则说明其潜在表示更有利于后续金融解释与截面决策。当前结果表明，在根仓库现有正式窗口内，Tucker 在这一点上更具优势。

进一步把排序结果转化为组合层指标后，可以观察到“排序相关性”和“组合表现”并不完全同步。从 formal 输出目录中的 \texttt{group\_returns\_overview.svg} 和 \texttt{drawdown\_overview.svg} 可以更直观看到这一点。在 baseline 结果中，HS300 上三类模型组合都取得了正累计收益，其中 Tucker 达到 2.11\%，明显高于 CP 的 0.87\% 和 PCA 的 0.74\%，且最大回撤仅为 $-0.03\%$；SZ50 上三类模型组合累计收益都为负，但 Tucker 的累计收益仅为 $-0.33\%$、最大回撤仅为 $-0.37\%$，明显优于 CP 的 $-3.00\%$ 累计收益和 $-3.63\%$ 最大回撤；ZZ500 上 PCA 与 Tucker 的累计收益分别为 $-1.49\%$ 和 $-1.43\%$，虽然都未转正，但已显著好于 CP 的 $-8.47\%$。换言之，单纯较高的 Rank IC 并不自动等价于更优的组合净值，这也是将组合结果与 Rank IC 联动分析的必要性所在。

在 extended 结果中，这种差异性更加明显。HS300 上扩展特征使三类模型的组合表现同步恶化，CP、PCA 和 Tucker 的累计收益分别降至 $-0.06\%$、$-2.70\%$ 和 $-4.26\%$，对应 Sharpe 也全部转为负值；SZ50 上扩展输入同样没有改善组合结果，三类模型累计收益进一步回落到 $-4.39\%$、$-2.44\%$ 和 $-3.51\%$。ZZ500 上则出现了相反现象，extended 输入显著改善了 CP 的组合表现，使其累计收益从 $-8.47\%$ 反转到 3.52\%，Sharpe 由 $-27.11$ 转为 16.99，而 Tucker 与 PCA 虽然回撤有所收敛，但累计收益仍未明显转正。这表明，扩展特征不仅会改变模型的重构与排序结果，也会改变它们在组合层的收益—风险权衡，因此在评价扩展输入时，必须同时观察净值曲线、回撤、换手率和行业/风格暴露，而不能只保留单一排序指标。

在进一步补充分位数组、多空组合、成本后净值和基准超额收益之后，可以更细致地判断“排序有效性是否真正能够转化为可配置收益”。例如，baseline HS300 上三类模型的 `top-bottom` 多空净值全部为正，其中 Tucker 与 CP 的多空累计净值分别达到 3.50\% 和 3.43\%，明显高于 PCA 的 2.57\%；但一旦计入交易成本，三类模型的成本后累计收益分别回落到 2.00\%、0.71\% 和 0.57\%，说明高换手模型的毛收益并不能无损落地。又如 baseline SZ50 上，Tucker 的多空组合累计净值达到 8.36\%，但相对指数基准的超额净值仍为负，说明其分位排序结构虽然存在，但在这一短窗口内尚不足以完全抵消市场基准波动。extended ZZ500 上，Tucker 与 PCA 的成本后净值分别仍保留 5.99\% 和 5.98\% 的正收益，且在修正为真实 ZZ500 基准后，相对基准的超额净值都在 3.86\% 左右，说明扩展输入在中盘样本上不仅改善了组合绝对收益，也部分增强了相对基准的配置价值。可见，加入分位数组、多空、成本和超额收益之后，组合层证据已经比单纯的 Top-N 净值曲线更接近真实投资决策场景。

为了把组合层结果与样本边界、行业暴露和风格暴露放到同一视角下比较，本文进一步生成了 \texttt{code/data/formal/reports/boundary\_portfolio/} 汇总表。该表同时列出各样本池的 Rank IC 最优模型、累计收益最优模型、Sharpe 最优模型，以及每个模型的 Top-N 收益、分组价差、多空净值、成本后净值、平均交易成本、超额净值、年化波动、平均换手、最大回撤和主要行业/风格暴露。汇总结果显示，Tucker 在 HS300、SZ50、ZZ500 和医药行业样本中往往兼具较好的排序和收益表现，但在全 A 样本中收益最优为 PCA，在大市值样本中 Rank IC 最优为 CP 而 Sharpe 最优为 Tucker。这说明模型优劣并不是单一指标可以决定的，必须同时观察排序相关性、组合收益、换手成本、超额表现和暴露集中度。

\section{扩展特征对照分析}

为了验证第一版扩展输入是否能够进入统一训练接口，本文进一步构造了 extended factor panel，并在 HS300、SZ50 与 ZZ500 三个样本池上分别运行 baseline/extended 两套 formal 配置。扩展版在原有四个主线因子之外，新加入市场级代理变量 \texttt{market\_return\_1d}、\texttt{market\_momentum\_5d}、\texttt{market\_momentum\_20d}、\texttt{market\_volatility\_20d}、\texttt{market\_drawdown\_20d}、\texttt{market\_amount\_change\_5d}、\texttt{market\_amount\_zscore\_20d}，财务 PIT 变量 \texttt{pit\_roe\_avg}、\texttt{pit\_np\_margin}、\texttt{pit\_gp\_margin}、\texttt{pit\_eps\_ttm}，以及业绩快报/业绩预告事件相关的 \texttt{perf\_express\_*}、\texttt{forecast\_*}、\texttt{event\_intensity\_score} 系列特征。对照结果表明，扩展输入的影响并不呈现单边改善，而是具有明显的样本池差异。

从已生成的对照结果看，HS300 上新增市场代理变量、宏观代理得分和事件字典衍生字段后，CP 的 Rank IC 从 0.0833 提升到 0.1771，但 Tucker 与 PCA 的 Rank IC 分别回落到 $-0.1428$ 和 $-0.1405$，且三类模型的组合 Sharpe 全部转负，说明扩展输入与模型结构之间仍存在显著相容性差异。SZ50 上，虽然 CP 的 MSE 从 0.9898 降至 0.7349，但三类模型的 Rank IC 全部下降，表明当前扩展特征对这一样本池的帮助仍然有限。ZZ500 上则出现了更有代表性的分化：CP 的 Rank IC 从 $-0.0779$ 提升到 0.0489，累计收益从 $-8.47\%$ 提升到 3.52\%；而 Tucker 与 PCA 的 Rank IC 仍接近于零，说明中盘样本对新增市场/事件信息更敏感，但这种敏感性并不对所有模型结构同时成立。

这一结果说明，第一版扩展输入已经成功完成“把市场代理变量、财务 PIT 与事件特征接入训练接口并形成可运行对照实验”的目标，但它还不能直接推出“扩展特征一定优于主线因子面板”的结论。更合理的解释是：市场级信息、PIT 财务特征和事件特征确实能够改变潜在结构，但这种改变是否真正提升模型质量，仍取决于样本池风格、模型结构以及特征筛选方式。因此，本文将这组结果视为后续继续扩展宏观变量、补做组合回测和进行样本边界分层实验的第一阶段证据。

\section{样本边界扩展探索}

为了进一步检验结论是否依赖于单一指数样本池，本文在当前提交版仓库中补充了样本边界扩展实验，包括全 A 活跃股票样本、三个行业分层样本以及大中小市值分层样本。与 HS300、SZ50、ZZ500 的正式基线相比，这批扩展实验更强调“不同样本边界下模型行为是否稳定”，首先在当前提交版短窗口口径下完成第一批 run，并进一步基于 long-window 因子面板对 2015、2020、2024 和 2026 四个代表性市场窗口做年度复核，重点观察 Rank IC、滚动稳定性和组合累计收益的相对变化。

从全 A 样本结果看，Tucker 依然是三类方法中最均衡的一类：其 Rank IC 为 0.0598，明显优于 CP 的 $-0.0997$ 和 PCA 的 $-0.0216$；同时，Tucker 的累计收益仅为 $-1.60\%$，也好于 CP 的 $-8.06\%$。这说明当样本边界从指数成分池扩展到更广泛的全 A 活跃股票后，三维结构保持带来的优势并没有完全消失，反而在更异质的样本中保留了更强的稳健性。

行业分层结果则表现出更强的结构差异。电子行业样本（C39）中，只有 CP 维持了 0.0571 的正向 Rank IC，而 Tucker 与 PCA 均转为负值，说明高波动成长行业中的截面信号更容易被特定模型结构放大或扭曲。医药行业样本（C27）中，Tucker 的 Rank IC 达到 0.2364，累计收益达到 27.81\%，明显好于 CP 和 PCA，说明在盈利质量和事件驱动更突出的行业里，三维结构压缩与排序的协同效应更强。专用设备行业样本（C35）中，PCA 的 Rank IC 最高，为 0.1392，而 Tucker 与 CP 相对较弱，说明不同行业边界下“最优模型”并不固定，样本风格会明显影响张量分解与二维压缩方法的相对优势。

市值分层结果进一步说明样本边界会改变排序与组合表现的对应关系。小市值样本中，CP 的 Rank IC 达到 0.2413，但组合累计收益仅为 $-0.58\%$，说明在高噪声、小样本场景下，排序相关性并不必然转化为可配置收益。中市值样本中，PCA 的 Rank IC 为 0.1416，累计收益达到 4.90\%，而 Tucker 仅取得轻微负收益，说明中盘样本更容易让二维压缩直接服务于组合层。大市值样本中，CP 的 Rank IC 最高，为 0.1451，但三类模型组合收益均为负，这再次表明“排序指标优于同业”与“组合净值最终为正”之间仍存在显著缺口。

在全年 long-window 复核中，这种“最优模型随边界和年份切换”的现象依然存在。根据 \texttt{code/data/formal/reports/long\_window\_portfolio/README.md}，全 A 样本在 2015 窗口中呈现“Rank IC 最优为 CP、收益最优为 PCA”的结构，2016 与 2017 两年则更偏向 Tucker，2020 窗口中收益最优转为 Tucker，2024 窗口中 Rank IC 最优再次回到 CP，2026 窗口则由 PCA 与 Tucker 分别在 Rank IC 和收益上占优。行业与市值分层同样没有表现出跨年份恒定不变的最优模型，例如 2021--2023 年的行业 C27 与 C39 更频繁出现 Tucker 最优，而 2018、2019、2025 等窗口里 PCA 或 CP 也会重新占优。总体而言，这些新增 run 已经把“样本边界扩展”从计划推进为真实产物，并证明结论不仅依赖样本边界，也依赖市场阶段。

从暴露角度看，行业分层样本的主要行业暴露自然集中在单一行业，而全 A 与指数样本更能暴露模型选择偏好。例如全 A 的 Tucker 组合主要暴露在医药制造、黑色金属和影视传媒等行业，并且风格暴露集中在质量因子；HS300 的 Tucker 组合则明显集中于货币金融服务，同时混合质量、价值和波动率因子。这个结果进一步提醒，模式发现不能只看潜在聚类是否稳定，也要检查组合层是否因为样本边界变化而集中到少数行业或单一风格。

\section{稳健性与局限性分析}

在滚动稳定性方面，HS300、SZ50 与 ZZ500 上的 Tucker 分别达到 0.9898、0.9923 和 0.9683，均显著高于 CP，也与 PCA 基本处于同一量级。稳定性的意义在于，模型在不同滚动窗口上学到的潜在结构具有较强一致性，这使得研究者能够更放心地解释因子重要性和时间状态变化，而不会因为模型在不同窗口下频繁漂移而失去可解释基础。对研究型实证而言，稳定性往往比单次高分更重要，因为它决定了结论是否具备可重复性。

\begin{figure}[htbp]
\centering
\begin{tikzpicture}
\begin{groupplot}[
  group style={group size=3 by 1, horizontal sep=1.3cm},
  width=0.28\textwidth,
  height=6.0cm,
  ybar,
  ymin=0.65,
  ymax=1.02,
  symbolic x coords={CP,Tucker,PCA},
  xtick=data,
  nodes near coords,
  nodes near coords align={vertical},
  every node near coord/.append style={font=\tiny},
  tick label style={font=\small},
  ylabel style={font=\small},
  title style={font=\small}
]
\nextgroupplot[title={HS300},ylabel={滚动稳定性}]
\addplot[fill=blue!55] coordinates {(CP,0.6974) (Tucker,0.9898) (PCA,0.9946)};
\nextgroupplot[title={SZ50}]
\addplot[fill=orange!70] coordinates {(CP,0.5320) (Tucker,0.9923) (PCA,0.9956)};
\nextgroupplot[title={ZZ500}]
\addplot[fill=teal!65] coordinates {(CP,0.8384) (Tucker,0.9683) (PCA,0.9575)};
\end{groupplot}
\end{tikzpicture}
\caption{三个正式样本池上的模型滚动稳定性对比图}
\label{fig:stability}
\end{figure}

当然，本文仍存在若干需要继续完善的方面。第一，在数据输入层面，当前正式训练输入已经在 extended 配置中接入市场状态代理、财务 PIT 指标和公告事件强度特征，并把字段字典、披露日和最早可用交易日规则写入扩展合同。但外部利率、宏观月度指标、分红、重大事项和公告文本等更宽口径来源尚未完成可用时点字段落盘，因此不能直接进入训练口径。这意味着当前模型仍主要是在“核心因子集合 + 已完成 PIT 校验的扩展输入”上验证张量分解方法，而对于更宽口径、多来源因子体系的适应性仍需在源表补齐后继续检验。

第二，在效果验证层面，本文已经补齐了 Top-N 组合的累计净值、回撤、换手率和行业/风格暴露等基础结果，并进一步加入了分位数组、多空组合、成本后净值、交易成本和相对样本池指数的超额收益，且已通过跨样本边界汇总表与 `2015-2026` 全年 long-window 组合汇总统一比较。该证据链已经能够从短窗口与全年窗口两个层面回答“若将排序结果转化为真实组合规则，会得到怎样的收益—风险画像”；当前仍需继续深化的是更细粒度的交易冲击刻画和风险归因模型，而不是再停留在“是否具备长窗口证据”这一层面。

第三，在样本边界层面，本文已经从 HS300、SZ50 与 ZZ500 三个指数样本池扩展到全 A 活跃股票、行业分层和市值分层实验，并为 2015--2026 年区间完成了按年度拆分的 long-window 真实复跑。新增结果说明，指数样本上的结论不能被机械外推：Tucker 在部分窗口和部分行业样本中表现较稳健，但电子、专用设备和不同市值分层样本的最优模型并不固定。虽然全年 long-window 结果已经补齐，但由于当前训练输入仍未纳入更广泛的宏观与事件源，这批扩展实验更适合被理解为“主线四因子体系在不同边界与市场阶段下的稳健性检验”，而不是所有候选输入都已完备后的最终全市场结论。

第四，在可视化与论证层面，本文已经补充解释方差、Rank IC、滚动稳定性、时间状态、因子热力图、股票潜在结构图、行业聚类交叉图和跨样本边界对比图，并将这些图表整理为答辩材料素材包，能够比单纯模型比较更完整地支撑模式发现叙事。本轮进一步补充了 2015--2026 各年度全 A 长窗口的时间状态切换图和因子重要性热力图，并按年度生成 `code/data/formal/reports/pattern_discovery_long_window_<year>/` 图组，用于替换此前“长窗口图待补”的占位说明。当前仍需继续完善的是把这些 long-window 图组与更宽输入边界实验联动，而不是继续停留在主线四因子输入上。

针对上述不足，本文给出如下完善计划。其一，在数据层面，继续按照“先补 PIT 安全、再扩特征输入”的顺序推进：已完成市场状态代理、财务 PIT 和公告事件特征的披露时点校验与落盘合同，后续只在外部宏观和更细粒度事件源表具备可用时点字段后继续接入。其二，在评估层面，现有链路已经覆盖 Top-N、分组收益、多空、成本后收益、超额收益、回撤、换手、波动、Sharpe 和暴露，后续重点转向长窗口复核和风险归因。其三，在样本层面，保留指数样本池实验作为稳定基线，同时使用全 A 股、行业分层和不同市值分层实验比较不同样本边界下潜在结构、最优秩和稳定性指标的差异。其四，在结果表达层面，继续扩展长窗口图形体系，并把现有模式发现图组、组合闭环报告和训练输入合同持续维护为答辩材料素材包。

\chapter{结论与展望}

\section{研究结论}

本文以“股票—因子—时间”三维张量为统一研究对象，围绕股票因子降维与模式发现问题，构建了正式数据链、张量建模框架和结构化实验输出体系。基于 HS300、SZ50 与 ZZ500 三个正式样本池在当前提交版输出窗口上的实验结果，本文得到以下主要结论。

第一，股票因子问题天然具有多模态结构，使用三阶张量表示比单纯二维矩阵表示更符合数据本质。张量表示不仅保留了股票截面、因子集合和时间演化的耦合关系，也为后续模式发现提供了天然结构基础。

第二，从数值重构角度看，PCA 作为二维压缩基线依然具有显著优势；但从结构保持与模式解释角度看，Tucker 分解在三个样本池上表现出更平衡的重构质量、排序能力与长期稳定性。尤其在 Rank IC 和滚动稳定性指标上，Tucker 相较于 CP 更稳定，也比单纯依赖 PCA 更能支撑结构解释与排序分析。

第三，张量分解不仅可以完成因子降维，还能够输出股票相似结构、因子共振关系和时间状态变化等模式信息。实验结果显示，质量、价值、动量和波动率等因子在潜在空间中承担了较重要的解释角色，不同时间阶段也存在显著的结构重排现象，说明张量方法适合从“模式发现”而不仅仅是“数值压缩”的视角理解股票因子系统。

\section{创新点与不足}

本文的创新主要体现在三个方面。其一，将股票因子降维、模式发现与正式数据工程链路放在统一框架下讨论，避免了“算法一套、数据一套、论文一套”的证据割裂。其二，在毕业论文层面同时比较 CP、Tucker 与 PCA 三类方法，不仅比较重构误差，还比较 Rank IC 与滚动稳定性，使结论更贴近金融研究需求。其三，将实验结果沉淀为结构化输出与图形化结果，使论文中的图表、结果表和仓库产物具有统一来源。

不足之处主要包括：虽然全 A、行业和市值分层实验已经补充到 2015--2026 全部 long-window 年度，并形成了 `long_window_portfolio` 汇总表与按年度组织的 long-window 模式发现图组，但当前训练输入仍主要停留在主线四因子与第一阶段 extended 输入上，尚未纳入更广泛的外部宏观、分红、重大事项和公告文本源表；当前实验已经形成 Top-N 组合净值、分位数组、多空、成本后净值、基准超额收益和跨样本边界统一汇总，但交易冲击和风险归因仍有继续细化空间；模式发现图表已经从基础指标扩展到股票潜在结构、行业聚类关系和全年 long-window 图组，但这些结果仍建立在当前输入边界之内。参考文献元数据已经按当前学校提交口径完成最终校验，后续只有在学院模板明确要求逐条展示 DOI 时才需要统一追加。

\section{后续工作展望}

后续工作可以从三个方向继续推进。第一，在数据层继续补齐外部宏观、分红、重大事项和公告文本的可用时点字段，在不破坏 PIT 安全的前提下构建更完整的多模态金融张量。第二，在方法层继续研究秩自动选择、稀疏约束、非负张量分解和动态张量分解，以提升模型解释性与适应性。第三，在应用层把当前结构化实验输出延伸到长窗口组合复核、回撤控制、风险归因和风格轮动监测，从而形成更完整的量化研究闭环。

\chapter*{参考文献}
\addcontentsline{toc}{chapter}{参考文献}
\begin{enumerate}[label={[\arabic*]}]
  \item Hitchcock F L. The Expression of a Tensor or a Polyadic as a Sum of Products[J]. Journal of Mathematical Physics, 1927, 6(1): 164-189.
  \item Tucker L R. Some Mathematical Notes on Three-Mode Factor Analysis[J]. Psychometrika, 1966, 31: 279-311.
  \item Harshman R A. Foundations of the PARAFAC Procedure: Models and Conditions for an Explanatory Multi-modal Factor Analysis[J]. UCLA Working Papers in Phonetics, 1970, 16: 1-84.
  \item Kolda T G, Bader B W. Tensor Decompositions and Applications[J]. SIAM Review, 2009, 51(3): 455-500.
  \item 胡胜龙. 张量分解的唯一性[J]. 运筹学学报(中英文), 2025, 29(03): 34-60.
  \item 夏虹, 张雅倩, 靳晓东, 陈彦萍, 高聪, 王忠民. 基于张量的方法及应用综述[J]. 计算机技术与发展, 2022, 32(06): 1-8.
  \item Rabanser S, Shchur O, Günnemann S. Introduction to Tensor Decompositions and their Applications in Machine Learning[EB/OL]. arXiv:1711.10781, 2017.
  \item 岳欣. 基于张量分解的股票波动分析[D]. 西南财经大学, 2022.
  \item 马少沛, 孙庆慧, 武雅萱, 等. 大数据下张量充分降维方法及其应用研究[J]. 统计研究, 2021, 38(02): 114-134.
  \item 曾亚丽. 基于张量模型的金融时间序列预测研究[D]. 武汉理工大学, 2017.
  \item Brandi G, Gramatica R, Di Matteo T. Unveil Stock Correlation via a New Tensor-based Decomposition Method[EB/OL]. arXiv:1911.06126, 2020.
  \item Spelta A. Financial Market Predictability with Tensor Decomposition and Links Forecast[J]. Applied Network Science, 2017, 2: 7.
  \item 黄宏运, 王梅, 朱家明. 基于多元回归分析的多因子选股模型[J]. 通化师范学院学报, 2016, 37(08): 44-46.
  \item 刘宇轩, 金伟泽, 袁亮. 多因子量化选股模型优化与实证研究——引入金融周期指标的分析[J]. 价格理论与实践, 2022(04): 141-145.
  \item Lettau M. High-Dimensional Factor Models with an Application to Mutual Fund Characteristics[R]. NBER Working Paper No.29833, 2022.
  \item Han Y F, Chen R, Zhang C H. Rank Determination in Tensor Factor Model[EB/OL]. arXiv:2011.07131, 2022.
  \item Wang D, Zheng Y, Lian H, Li G D. High-dimensional Vector Autoregressive Time Series Modeling via Tensor Decomposition[EB/OL]. arXiv:1909.06624, 2020.
  \item 李博, 孟志青, 朱爱花. 时态支持向量机模型在股票操纵模式发现上的研究[J]. 系统科学与数学, 2023, 43(02): 356-378.
  \item 徐常青, 祁力群. 张量CP分解、半正定张量和范德蒙张量[J]. 苏州科技学院学报(自然科学版), 2016, 33(02): 1-7+82.
\end{enumerate}

\appendix
\chapter{实验参数与补充说明}

本文正式实验以 HS300、SZ50 与 ZZ500 为样本池。当前根仓库中已提交的正式结果窗口为 2026 年 3 月 2 日至 2026 年 3 月 30 日，模型层固定比较 CP、Tucker 与 PCA 三类方法，评价指标包括 MSE、解释方差、Rank IC、IR 与滚动稳定性。实验结果统一由项目输出目录中的结构化文件生成，确保论文中的图表、表格和文本结论具有一致证据来源。
